\documentclass[aps,pre,preprint,superscriptaddress,nofootinbib]{revtex4-2}
\usepackage{amsmath,amssymb}
\usepackage{hyperref}
\begin{document}
\title{Sharp lower bounds on current cumulant ratios in Markov rings with an irreversible bond}
\author{Xiaojian Zeng}
\email{xiaojianzeng@foxmail.com}
\affiliation{PBC School of Finance, Tsinghua University, 100190, China}
\date{\today}
\begin{abstract}
For a finite continuous-time Markov ring with positive forward rates, nonnegative reverse rates, and at least one irreversible bond, we prove the sharp lower bounds
$c_3/c_1\geq-1/16$ and $c_4/c_2\geq-(1+\sqrt{5})/10$ for the long-time cumulant rates of a signed unit cycle count.
The constants were previously proposed for unicyclic currents.
An irreversible bond makes the count regenerative; the established birth-and-death passage-time representation reduces the problem to inequalities for arbitrary positive exponential-stage weights.
The proof allows unrestricted backward excursions on the remaining bonds and does not assume a slow-link extremal topology.
Sharpness is approached by directed rings of increasing size.
The conjecture for rings with all reverse rates strictly positive remains unresolved.
\end{abstract}
\maketitle

\section{Introduction}
Higher current cumulants can constrain the microscopic dynamics compatible with a transport measurement. Ptaszy\'nski proposed the lower constants $-1/16$ and $-(1+\sqrt{5})/10$ for the ratios of third to first and fourth to second steady-state current cumulants in unicyclic Markov networks \cite{Ptaszynski2022}. The directed slow-link family yields those constants in the large-network limit; selecting that family as globally extremal relies on numerical evidence in the original treatment.

Here we prove both inequalities for every ring containing an irreversible bond, including rings with arbitrary backward rates elsewhere. The contribution is a proof without an extremal-topology ansatz in this boundary class. The constants, renewal relations, and passage-time spectral representation are established prior results. The argument does not establish the conjecture for strictly positive reverse rates.

The irreversible-bond model also occurs in statistical kinetics \cite{BaratoSeifert2015}. Bounds on standardized cycle-completion-time or first-passage-time moments \cite{BaratoSeifert2015,WamplerBarato2021} concern different observables from the current cumulant ratios used here. The distinction matters because passing from a waiting time to its inverse counting process couples several waiting-time cumulants.

\section{Model and result}
Let $N\geq3$ and label ring states modulo $N$. The forward transition $i\to i+1$ has finite rate $u_i>0$, and its reverse has finite rate $w_i\geq0$. Assume $w_i=0$ on at least one bond. There are no other transitions and the rates are time independent. With column probability vectors the generator is
\begin{equation}
 L_{i+1,i}=u_i,\qquad L_{i,i+1}=w_i,\qquad
 L_{i,i}=-(u_i+w_{i-1}).
 \label{eq:generator}
\end{equation}
Thus every column sums to zero. The diagonal follows directly from outgoing escapes; all counting-field factors below modify off-diagonal jump rates only.

For any fixed bond, let $Q_t$ be its net forward crossing count, with increments $+1$ and $-1$. Define
\begin{equation}
 \psi(z)=\lim_{t\to\infty}\frac{1}{t}
 \log\mathbb E[e^{zQ_t}],\qquad c_n=\psi^{(n)}(0),\qquad
 S=\frac{c_3}{c_1},\quad K=\frac{c_4}{c_2}.
 \label{eq:observable}
\end{equation}
These are cumulant ratios, not standardized central moments.
Rates and $c_n$ have units of inverse time; $Q_t,z,S,K$ are dimensionless.
Positive forward rates ensure irreducibility. The finite tilted generator has an analytic simple eigenvalue branch near $z=0$, which defines $\psi$ there for any initial law.

For this entire class,
\begin{equation}
 \boxed{S\geq-\frac{1}{16},\qquad
 K\geq-\frac{1+\sqrt{5}}{10}.}
 \label{eq:bounds}
\end{equation}
Both denominators are strictly positive. Both constants are sharp over unbounded finite $N$. The proof and the limiting families follow below.

\section{Renewal reduction}
Rotate the labels so $w_0=0$ and count the transition $0\to1$. Each crossing returns the system to state 1. Before the next crossing, unroll the states as $1,2,\ldots,N-1,0$, replacing the next $0\to1$ jump by absorption. This is a birth-and-death chain with $N$ transient states, reflecting entrance, and absorbing exit. The strong Markov property gives independent identically distributed intercrossing times $\tau$.

Absorption has an exponential tail: from every transient state, a sequence of at most $N$ forward jumps within a fixed time has a uniformly positive probability. A finite initial delay therefore leaves the long-time cumulant rates unchanged. Net counts on two bonds differ by the change in occupancy of the intervening arc, which is bounded. Their real counting-field generating functions differ by bounded factors, so their scaled logarithms, and hence all derivatives of the analytic eigenvalue branch, coincide.

Let $T$ be the transient subgenerator. The birth-and-death passage-time theorem \cite{Fill2009}, including its extension to vanishing death rates, gives
\begin{equation}
 \widehat F(a):=\mathbb E[e^{-a\tau}]
 =\prod_{j=1}^N\frac{\lambda_j}{\lambda_j+a},
 \label{eq:transform}
\end{equation}
where $\lambda_j$ are the positive eigenvalues of $-T$, with multiplicity.
For completeness, the endpoint cofactor of the tridiagonal resolvent gives $\widehat F(a)=(\prod_i u_i)/\det(aI-T)$.
Eventual absorption fixes $\det(-T)=\prod_i u_i$.
Each communicating interval of $T$ is diagonally symmetrizable; zero reverse rates make the full matrix block triangular with these interval blocks. Stability from the exponential absorption tail makes all $\lambda_j$ real and strictly positive. Repeated eigenvalues pose no difficulty. Equation~(\ref{eq:transform}) is the transform of a sum of independent exponential stages.

Write
\begin{equation}
 x_j=\lambda_j^{-1},\quad M=\sum_jx_j,\quad
 p_j=\frac{x_j}{M},\qquad
 q=\sum_jp_j^2,\quad r=\sum_jp_j^3,\quad h=\sum_jp_j^4.
 \label{eq:weights}
\end{equation}
All $p_j$ are positive and sum to one. The waiting-time cumulants are $\kappa_n=(n-1)!\sum_jx_j^n$.
The renewal relation $\log\widehat F(\psi(z))=-z$ yields the familiar inversion formulas \cite{Ptaszynski2022}
\begin{align}
 c_1&=\kappa_1^{-1},&
 c_2&=\frac{\kappa_2}{\kappa_1^3},\nonumber\\
 c_3&=\frac{3\kappa_2^2}{\kappa_1^5}
       -\frac{\kappa_3}{\kappa_1^4},\nonumber\\
 c_4&=\frac{15\kappa_2^3}{\kappa_1^7}
       -\frac{10\kappa_2\kappa_3}{\kappa_1^6}
       +\frac{\kappa_4}{\kappa_1^5}.
\end{align}
Consequently $c_1=1/M>0$, $c_2=q/M>0$, and
\begin{equation}
 S=3q^2-2r,\qquad K=15q^2-20r+\frac{6h}{q}.
 \label{eq:ratios}
\end{equation}
No approximation in ring size or rate ratios has entered this reduction.

\section{Moment inequalities}
\subsection{Third cumulant ratio}
Finite-vector norm monotonicity gives $r\leq q^{3/2}$, and $0<q\leq1$. With $x=\sqrt q$,
\begin{equation}
 S\geq3x^4-2x^3\geq-\frac1{16},
\end{equation}
since
\begin{equation}
 3x^4-2x^3+\frac1{16}
 =\frac{(2x-1)^2(12x^2+4x+1)}{16}\geq0.
\end{equation}
This argument holds for every positive stage-weight vector.

\subsection{Fourth cumulant ratio}
Cauchy--Schwarz gives $r^2\leq qh$. Set $y=r/q$, so $0<y\leq\sqrt q$. Equation~(\ref{eq:ratios}) implies
\begin{equation}
 K\geq15q^2-20qy+6y^2.
 \label{eq:kquadratic}
\end{equation}
When $q\leq9/25$, completing the square gives
\begin{equation}
 K\geq-\frac53q^2\geq-\frac{27}{125}
 >-\frac{1+\sqrt5}{10}.
\end{equation}
When $q\geq9/25$, the derivative of the right-hand side of Eq.~(\ref{eq:kquadratic}) with respect to $y$ is
$12y-20q\leq12\sqrt q-20q\leq0$.
Thus replacing $y$ by $\sqrt q$ gives a lower bound:
\begin{equation}
 K\geq f(x):=15x^4-20x^3+6x^2,\qquad \frac35\leq x\leq1.
\end{equation}
Let $p_*=(5+\sqrt5)/10$. The exact identity
\begin{align}
 f(x)+\frac{1+\sqrt5}{10}
 &=(x-p_*)^2
 \left[15x^2+(-5+3\sqrt5)x+\frac{\sqrt5-1}{2}\right]
 \geq0
\end{align}
proves the second bound. The bracket coefficients are positive for $x\geq0$.
The split at $q=9/25$ is necessary to justify the substitution in Eq.~(\ref{eq:kquadratic}).

\section{Sharpness}
Choose a time scale $k>0$ and directed rings ($w_i=0$ for all $i$). To approach the third-ratio bound take
\begin{equation}
 u_0=\frac{k}{N-1},\qquad u_i=k\quad(i>0).
\end{equation}
The normalized stage weights are $p_0=1/2$ and $p_i=1/[2(N-1)]$ for $i>0$. Hence $q\to1/4$, $r\to1/8$, and $S\to-1/16$.

For the fourth-ratio bound take instead
\begin{equation}
 u_0=\frac{k(1-p_*)}{p_*(N-1)},\qquad u_i=k\quad(i>0).
\end{equation}
Then $p_0=p_*$ and $p_i=(1-p_*)/(N-1)$. As $N\to\infty$, $(q,r,h)\to(p_*^2,p_*^3,p_*^4)$, so $K\to f(p_*)=-(1+\sqrt5)/10$.
The slow-rate prefactor is $(1-p_*)/p_*=(3-\sqrt5)/2$.
These families agree with the earlier limiting construction \cite{Ptaszynski2022}. Every member has finite positive forward rates; sharpness does not assert equality or characterize the optimizer at a fixed finite $N$.

\section{Checks and physical scope}
The analytic argument establishes the inequalities for arbitrary finite ring size. Exact finite-case calculations support the transcription and counting conventions: the accompanying standard-library Python code compares backward absorption-moment inversion with normalized tilted-generator differentiation on all counted bonds of 16 rings at $N=3,4,6,8$. It includes directed rings, one or several irreversible bonds, and strong internal backward excursions; it makes 84 exact four-cumulant comparisons. It also checks the scalar polynomial identities in rational arithmetic and in $\mathbb Q(\sqrt5)$. These checks supplement the proof rather than extend its domain.

The bounds constrain transport models containing an irreversible step even when the other steps have substantial backtracking. A measured long-time ratio below either bound excludes this class provided unit cycle counting, time-independent rates, and the single-ring topology apply. It does not identify which premise fails.
An exactly irreversible bond corresponds to infinite cycle affinity under a local-detailed-balance interpretation. No finite-affinity entropy-production bound is claimed. Additional channels, finite-time ratios, time-dependent rates, and differently normalized observables require separate analyses. In particular, continuity from a zero reverse rate proves no bound throughout the strictly positive interior.

\subsection*{Research tools and data availability}
AI assistance through OpenAI Codex was used for literature comparison, analytic reasoning, code generation, exact calculation checks, and drafting. This writing session used a GPT-6-based agent; exact model versions for earlier sessions were not recorded and remain to be confirmed. The work was directed toward a fixed Markov-ring claim and checked using explicit algebra and separate exact generator and passage-time calculations, including a fresh AI-assisted local scientific audit. Such checks do not replace human scientific judgment. Human author review is pending for this local draft.
The accompanying reproducibility files specify the deterministic test inputs, algorithms, environment, and execution instructions. There are no experimental data.

\begin{acknowledgments}
OpenAI Codex also assisted with manuscript organization and wording. The scope and accuracy of this disclosure require author confirmation before submission.
\end{acknowledgments}
\bibliographystyle{apsrev4-2}
\bibliography{references}
\end{document}
